\begin{table*}[tbp]
\centering
\footnotesize
\setlength{\tabcolsep}{4pt}
\caption{Effect of an error in the secondary-path model added to the truncated model (600 runs per algorithm and error level)}
\label{tab:atsip-7}
\par\smallskip
\begin{tabular}{lllllllll}
\toprule
 & \multicolumn{4}{l}{Diverged runs} & \multicolumn{4}{l}{Median recovery time [s]} \\
 & & Without & CohVSS- & & \multicolumn{2}{l}{Proposed} & \multicolumn{2}{l}{\makecell[l]{CohVSS-\\FxNLMS}} \\
\makecell[l]{Added\\error} & Proposed & trials & FxNLMS & C-FXLMS & Close & Far & Close & Far \\
\midrule
None & 0 & 0 & 0 & 0 & 0.36 & 0.40 & 0.98 & 1.26 \\
$-$40\,dB & 0 & 0 & 0 & 0 & 0.36 & 0.41 & 0.98 & 1.26 \\
$-$30\,dB & 9 & 6 & 3 & 3 & 0.37 & 0.43 & 0.98 & 1.26 \\
$-$20\,dB & 54 & 51 & 21 & 54 & 0.44 & 0.51 & 1.03 & 1.33 \\
\bottomrule
\end{tabular}
\par\medskip
\begin{minipage}{0.98\textwidth}
\scriptsize
The model already contains a truncation error of about $-$24\,dB. The 600 runs consist of the no-change, close-source, and far-source conditions. The recovery time is computed over the runs without divergence.\par
\end{minipage}
\end{table*}
